\documentclass[10pt,a4paper]{article}
\usepackage{amsmath,amssymb,amsthm,graphicx}
% Minimal retained preamble for lane-p2 cards taken from
% sources/ahl-16-fixed-points-2019/editorial-preamble.tex
% (doi:10.5802/ahl.16, Annales Henri Lebesgue 2 (2019) 149-185, CC BY 4.0).
%
% Copied verbatim from the editorial preamble: the package list, the
% theorem declarations under the shared counter and the macro block, so that
% every macro and environment of the retained spans is available.
%
% Omitted and why: the class line (the assembler supplies the class); the
% a4paper geometry line and the fontspec block with its explicit lmroman10
% files (encodings and fonts fight the forced xelatex engine, and the
% assembler supplies the page setup).
\usepackage{amsmath,amssymb,amsthm,mathtools,graphicx,xcolor,hyperref,xurl,xspace,ifthen}
\newtheorem{theo}{Theorem}[section]
\newtheorem{lemm}[theo]{Lemma}
\newtheorem{theorem}[theo]{Theorem}
\newtheorem*{remas*}{Remarks}
\providecommand{\og}{«\nobreakspace}
\providecommand{\fg}{\nobreakspace»}
\newtheorem{lem}[theo]{Lemma}
\newtheorem{lemma}[theo]{Lemma}
\newtheorem{prop}[theo]{Proposition}
\newtheorem{coro}[theo]{Corollary}
\newtheorem{corol}[theo]{Corollary}
\newtheorem{conj}[theo]{Conjecture}
\newtheorem{Proposition}[theo]{Proposition}
\newtheorem{theoreme}[theo]{Theorem}
\theoremstyle{definition}
\newtheorem{defi}[theo]{Definition}
\newtheorem{exam}[theo]{Example}
\newtheorem{exem}[theo]{Example}
\newtheorem{rema}[theo]{Remark}
\newtheorem{nota}[theo]{Notation}
\newtheorem*{rema*}{Remark}
\newtheorem*{defi*}{Definition}
\newtheorem*{theo*}{Theorem}
\newtheorem*{lemm*}{Lemma}
\newtheorem*{prop*}{Proposition}
\newcommand{\enoncename}{}
\newtheorem{corpusenonce}[theo]{\enoncename}
\newenvironment{enonce}[2][]{\renewcommand{\enoncename}{#2}\begin{corpusenonce}}{\end{corpusenonce}}
\newenvironment{enonce*}[2][]{\par\medskip\noindent\textbf{#2.}\itshape}{\par\medskip}
\newenvironment{proofc}[1][\proofname]{\begin{proof}[#1]}{\end{proof}}
\newcommand{\Mc}{\mathcal{M}}
\newcommand{\pointir}{.\ ---\ }
\providecommand{\mathof}[1]{\mathrm{#1}}

\numberwithin{equation}{section}
\numberwithin{figure}{section}
\numberwithin{table}{section}
\def\endappendix{}
%~Mouliné par MaN_auto v.0.21.4 2019-04-30 09:02:09


\usepackage[all]{xy}

\hyphenation{auto-morphism
every
Roelcke
peri-od-ic
peri-od-icy
rep-re-sen-ta-tion
}








\newcommand{\A}{\mathcal A}
\newcommand{\N}{\mathbb{N}}
\newcommand{\Z}{\mathbb{Z}}
\newcommand{\Q}{\mathbb{Q}}
\newcommand{\K}{\mathcal{K}}
\newcommand{\CC}{\mathbb{C}}

\newcommand*{\mA}{\mathbf{A}}
\newcommand*{\mB}{\mathbf{B}}
\newcommand*{\mC}{\mathbf{C}}
\newcommand*{\mD}{\mathbf{D}}
\newcommand*{\mE}{\mathbf{E}}
\newcommand*{\mF}{\mathbf{F}}
\newcommand*{\mH}{\mathbf{H}}
\newcommand*{\mZ}{\mathbf{Z}}
%\newcommand*{\mA}{\mathbf{A}}
%\newcommand*{\mA}{\mathbf{A}}

\newcommand{\restrict}[2]{#1 \upharpoonright #2}

\newcommand{\e}{\varepsilon}

\DeclareMathOperator{\Aut}{Aut}
\DeclareMathOperator{\Stab}{Stab}
\DeclareMathOperator{\Age}{Age}

\newcommand{\actson}{\curvearrowright}

\newcommand{\RUCB}{\mathrm{RUC}_{b}}
\newcommand{\ROB}{\mathrm{Ro}_{b}}
\newcommand{\WAP}{\mathrm{WAP}}
\newcommand{\AP}{\mathrm{AP}}

\graphicspath{{./figures/}}

\newcommand*{\mk}{\mkern -1mu}
\newcommand*{\Mk}{\mkern -2mu}
\newcommand*{\mK}{\mkern 1mu}
\newcommand*{\MK}{\mkern 2mu}

\hypersetup{urlcolor=purple, linkcolor=blue, citecolor=red}

\newcommand*{\romanenumi}{\renewcommand*{\theenumi}{\roman{enumi}}}
\newcommand*{\Romanenumi}{\renewcommand*{\theenumi}{\Roman{enumi}}}
\newcommand*{\alphenumi}{\renewcommand*{\theenumi}{\alph{enumi}}}
\newcommand*{\Alphenumi}{\renewcommand*{\theenumi}{\Alph{enumi}}}
\let\oldexists\exists
\renewcommand*{\exists}{\mathrel{\oldexists}}
\let\oldforall\forall
\renewcommand*{\forall}{\mathrel{\oldforall}}


\title{Two equivalent formulations of the almost-periodicity criterion for finite colorings of a Fraisse structure}
\author{Lionel Nguyen Van Thé}
\date{doi:10.5802/ahl.16; CC BY 4.0}
\begin{document}
\maketitle
\noindent\textit{Editorial note.} The counted unit is the article's own Proposition prop:aOSRamsey at source lines 318-325 together with its proof at source lines 327-338, a straight adjacent pairing in the source file. The source is the publisher version of the article in Annales Henri Lebesgue 2 (2019), pages 149-185, DOI 10.5802/ahl.16, released under CC BY 4.0.\par\medskip
\noindent\textit{Editorial note.} The licence evidence for this card is the inherited publisher-version record kept with the source: LICENSE\_EVIDENCE.json names the publisher version, the licence CC BY 4.0 and its URL, quotes the notice on the closing page of the publisher PDF, and records the SHA-256 of that PDF, while inherited-license.txt restates author, work, version, DOI and licence. The source's own licence-observation.json records that the landing page displays no licence line and that the reuse basis is the inherited evidence; this lane verified the first article of the family and found no contradiction with the journal-level CC BY 4.0 statement, which is the check the Lead required for the AHL family. No acceptance record is written anywhere.\par\medskip
\noindent\textit{Editorial note.} The retained preamble is the source's editorial preamble with three omissions only: the class line, which the assembler supplies; the page geometry; and the fontspec block with its explicit font files, which would fight the forced xelatex engine. The theorem declarations under the shared counter and the whole macro block are kept verbatim, so that the macros of the retained spans, including the boldface structure letter, the age symbol, the stabiliser and the restriction notation, are all available.\par\medskip
\noindent\textit{Editorial note.} Neither the retained statement nor the retained proof contains a citation or a reference to another labelled result of the article, so this card carries no bibliography entry and no navigation pointer.\par\medskip
\noindent\textit{Editorial note.} The statements in the census package were used as a starting point only: the statement span, the proof span and the absence of citations were re-read from the source before the card was written, as the census itself asks.\par\medskip
\noindent\textit{Editorial note.} rqual sv-mismatch batch fix 2026-10-14 (Chief ruling: re-assemble unpublished only): source\_version aligned verbatim to the source dir license-observation.json sv (package evidence authority); revision bumped for re-assembly. Recipe from lane-s3.\par\medskip
\noindent\textit{Editorial note.} Revision r003 of this card aligns the licence evidence pointer with the source observation file kept with the source, sources/ahl-16-fixed-points-2019/license-observation.json. The counted statement, the proof spans and the navigation pointers are unchanged; only the pointer, and where it differed the recorded version string, moved.\par\medskip
\noindent\textit{Allowed standard prerequisites (editorial).}\begin{itemize}
\item Fraisse structures, their ages, and the quotients of the automorphism group. The paper works with a Fraisse structure based on the natural numbers, its age, that is the class of its finite substructures, and its automorphism group; the finite sets of embeddings of a finite substructure into the structure are written as binomial symbols, and the quotient of the automorphism group by the stabiliser of a finite substructure is identified with the set of embeddings of that substructure, whose image under an automorphism is again such a set. A colouring of such a set of embeddings is a map into a finite set of colours, and being constant up to a positive number on a set of embeddings means that its values on that set vary by at most that number. The author's own notation is retained inside the counted statement and proof.. The counted proposition is the combinatorial criterion that the paper uses to pass between statements about finite colorings and statements about the orbits of the automorphism group, and it is invoked by the Ramsey-type results of the same section.
\end{itemize}\medskip
\begin{prop}\label{prop:aOSRamsey}
Let $\mA \in \Age(\mF)$, $\mathcal C$ be a finite set of finite colorings of $\binom{\mF}{\mA}$, $\e>0$. TFAE:
\begin{enumerate}
\item For every finite $H\subset G$, there exists $g\in G$ so that every $\chi \in \mathcal C$ is constant up to $\varepsilon$ on $Hg$.

\item For every $\mB \in \Age(\mF)$, there exists $b\in \binom{\mF}{\mB}$ such that every $\chi\in \mathcal C$ is constant up to $\e$ on $\binom{b(\mB)}{\mA}$.
\end{enumerate}
\end{prop}

\begin{proof}
The proof hinges on the following observation. Let $\mA$ be a finite substructure of $\mF$ and $H$ be a finite subset of $\Aut(\mF)$. For $h\in \Aut(\mF)$, recall that $\bar h$ denotes the equivalence class of $h$ in the quotient $\Stab(\mA)\backslash \Aut(\mF)$. As we have seen, $\bar h$ can be thought of as the restriction $\restrict{h^{-1}}{A}$, so $\overline H:=\{ \bar h:h\in H\}$ can be seen as a finite set of embeddings of $\mA$ into $\mF$. As such, it is contained in some set of the form $\binom{\mB}{\mA}$ for some finite substructure $\mB$ of $\mF$. Next, if $g$ is fixed in $G$, we have:
\[
\overline{Hg} = \{ \overline{hg}:h\in H\}=\{ \restrict{g^{-1}\circ h^{-1}}{A} : h\in H\} \subset \binom{g^{-1}(\mB)}{\mA}.
\]

Conversely, if $\mB \subset \mF$ is a finite substructure, then there is $H\subset \Aut(\mF)$ finite so that $\binom{\mB}{\mA}\subset \overline H$, and if $g\in \Aut(\mF)$, then $\binom{g^{-1}(\mB)}{\mA} \subset \overline{Hg}$.

We now go on with the proof. Assume that for every finite $H\subset G$, there exists $g\in G$ so that every $\chi \in \mathcal C$ is constant up to $\varepsilon$ on $Hg$. Let $H\subset \Aut(\mF)$ be a finite set so that $\binom{\mB}{\mA} \subset \overline H$. Find $g\in \Aut(\mF)$ such that every $\chi\in \mathcal C$ is constant up to $\varepsilon$ on $Hg$. Then every $\chi\in \mathcal C$ is constant up to $\e$ on $Hg$, and hence on $\overline{Hg}\supset \binom{g^{-1}(\mB)}{\mA}$. Therefore, it suffices to set $b=\restrict{g^{-1}}{B}$.

Conversely, fix $H\subset \Aut(\mF)$ finite. Let $\mB$ be a finite substructure of $\mF$ so that $\overline{H} \subset \binom{\mB}{\mA}$. By hypothesis, find $b\in \binom{\mF}{\mB}$ such that every $\chi\in \mathcal C$ is constant up to $\e$ on $\binom{b(\mB)}{\mA}$. Take $g\in \Aut(\mF)$ such that $g^{-1}$ extends $b$. Then $Hg\subset \overline{Hg}\subset \binom{g^{-1}(\mB)}{\mA}$ and every $\chi\in \mathcal C$ is constant up to $\varepsilon$ on $Hg$.
\end{proof}

\end{document}
